\documentclass[11pt]{article}
\usepackage[margin=1in]{geometry}
\usepackage{amsmath,amssymb,amsthm,mathtools}
\usepackage{graphicx}
\usepackage{microtype}
\usepackage[colorlinks=true,linkcolor=blue,citecolor=blue,urlcolor=blue]{hyperref}
\usepackage{enumitem}

\newtheorem{theorem}{Theorem}
\newtheorem{corollary}[theorem]{Corollary}
\newtheorem{remark}[theorem]{Remark}
\newcommand{\R}{\mathbb R}
\newcommand{\eps}{\varepsilon}
\newcommand{\Dist}{\mathsf D}

\makeatletter
\renewcommand{\@seccntformat}[1]{\csname the#1\endcsname\quad}
\makeatother

\title{Cube Transfer for Randomized Dvoretzky in $\ell_p^n$\\
\large An $\eps$-dependent logarithmic regime and a uniform threshold improvement}
\author{Candidate partial result for arXiv:1510.07284\\
Run \texttt{fa\_banach\_001}, lane 0}
\date{9 August 2026}

\begin{document}
\maketitle

\noindent\fbox{\parbox{0.94\textwidth}{
\textbf{Status.} Candidate strong partial result, likely valid, pending
specialist human review.  The sharp $\eps/\log(1/\eps)$ upper bound is proved
when $p\ge C\eps^{-1}\log n$, and the source's uniform threshold
$p>(\log n)^2$ is improved to
$p\ge C(\log n)^2/\log\log n$ for nontrivial dimensions $k\ge2$.
The full logarithmic-$p$, very-small-distortion wedge remains open.}}

\section{Source question}

Paouris, Valettas and Zinn study the largest dimension of a Haar-random
subspace of $\ell_p^n$ on which the norm is $(1+\eps)$-spherical
\cite{PVZ17}.  In Section 6, item 5, they ask whether the factor
$\log(1/\eps)$ in the denominator of
\[
   k\lesssim \frac{\eps\log n}{\log(1/\eps)}
\]
is necessary throughout the range $p\gtrsim\log n$.  Their Proposition 6.2
proves this upper bound for $p>(\log n)^2$.

\begin{figure}[h]
\centering
\includegraphics[width=\textwidth]{figures/open_problem_crop.png}
\caption{The open question and Proposition 6.2 in arXiv:1510.07284,
PDF page 39.}
\end{figure}

The present packet improves the known parameter range but does not close the
whole question.

\section{\hspace{0.25em}Definitions and external input}

For a subspace $E\subset\R^n$ define its $\ell_q$ spherical distortion by
\[
 \Dist_q(E):=
 \frac{\displaystyle\sup_{x\in E\cap S^{n-1}}\|x\|_q}
      {\displaystyle\inf_{x\in E\cap S^{n-1}}\|x\|_q}.
\]
Thus $E$ is $(1+\eta)$-spherical in $\ell_q^n$ precisely when
$\Dist_q(E)\le1+\eta$.

We use the sharp randomized cube theorem of Tikhomirov \cite{Tik14}: there is
an absolute constant $C_T$ such that, if $k>1$, $0<\eta<1/2$, and a majority
of Haar-random $k$-subspaces of $\R^n$ satisfy
$\Dist_\infty(E)\le1+\eta$, then
\begin{equation}\label{eq:cube}
 k\le C_T\frac{\eta\log n}{\log(1/\eta)}.
\end{equation}
The probability threshold $3/4$ used below is stronger than ``a majority.''

\section{Proof intuition}

The key idea is to transfer a good $\ell_p$ section to the cube without
changing the random subspace.  Globally on $\R^n$ one has
$\|x\|_\infty\le\|x\|_p\le n^{1/p}\|x\|_\infty$.  On a fixed Euclidean
sphere this implies
\[
   \Dist_\infty(E)\le n^{1/p}\Dist_p(E).
\]
Hence a $(1+\eps)$-spherical $\ell_p$ section is a
$(1+\delta)$-spherical cube section with
$1+\delta=(1+\eps)n^{1/p}$.  Tikhomirov's cube obstruction then applies
verbatim.  When $(\log n)/p$ is at most a constant multiple of $\eps$, this
has the desired sharp dependence on $\eps$.  Uniformly in $\eps$, the same
bound excludes every $k\ge2$ once
$p\gtrsim(\log n)^2/\log\log n$ and $\eps$ lies below the transfer error.

\section{Main transfer theorem}

\begin{theorem}[Cube-transfer upper bound]\label{thm:transfer}
Let $n>1$, $1\le p<\infty$, $k>1$, and $0<\eps<1$.  Set
\begin{equation}\label{eq:delta}
  \delta:=(1+\eps)n^{1/p}-1.
\end{equation}
Suppose $0<\delta<1/2$ and that a Haar-random $k$-dimensional subspace $E$
of $\R^n$ satisfies $\Dist_p(E)\le1+\eps$ with probability greater than
$3/4$.  Then
\begin{equation}\label{eq:main}
  k\le C_T\frac{\delta\log n}{\log(1/\delta)},
\end{equation}
where $C_T$ is the absolute constant in the randomized cube theorem.
\end{theorem}

\begin{proof}
For every $x\in\R^n$,
\begin{equation}\label{eq:normcompare}
 \|x\|_\infty\le\|x\|_p\le n^{1/p}\|x\|_\infty.
\end{equation}
Fix a subspace $E$.  The first inequality gives
\[
 \sup_{E\cap S^{n-1}}\|x\|_\infty
 \le \sup_{E\cap S^{n-1}}\|x\|_p,
\]
whereas the second gives
\[
 \inf_{E\cap S^{n-1}}\|x\|_\infty
 \ge n^{-1/p}\inf_{E\cap S^{n-1}}\|x\|_p.
\]
Dividing yields
\begin{equation}\label{eq:distcompare}
   \Dist_\infty(E)\le n^{1/p}\Dist_p(E).
\end{equation}
Consequently, every $E$ with $\Dist_p(E)\le1+\eps$ satisfies
\[
 \Dist_\infty(E)\le n^{1/p}(1+\eps)=1+\delta.
\]
Thus more than $3/4$ of the Haar-random $k$-subspaces are
$(1+\delta)$-spherical in the cube.  Equation \eqref{eq:main} is exactly
\eqref{eq:cube} with $\eta=\delta$.
\end{proof}

\section{Consequences for the source question}

\begin{corollary}[$\eps$-dependent logarithmic regime]\label{cor:eps}
There is an absolute constant $C$ such that the following holds.  Let
$0<\eps<1/3$ and
\[
      p\ge 4\eps^{-1}\log n.
\]
If a Haar-random $k$-subspace of $\ell_p^n$ is $(1+\eps)$-spherical with
probability greater than $3/4$, then
\begin{equation}\label{eq:sharp}
      k\le C\frac{\eps\log n}{\log(1/\eps)}.
\end{equation}
\end{corollary}

\begin{proof}
Put $s=(\log n)/p$, so $s\le\eps/4$ and
$\delta=(1+\eps)e^s-1$.  For $0<\eps<1/3$, elementary estimates on the
exponential give $\delta<1/2$ and $\delta\le(3/2)\eps$.  On this interval,
\[
 \frac{\delta}{\log(1/\delta)}
 \le C_0\frac{\eps}{\log(1/\eps)}
\]
for an absolute $C_0$.  Theorem \ref{thm:transfer} proves
\eqref{eq:sharp}.
\end{proof}

In particular, for every fixed distortion $\eps>0$, the logarithmic loss is
necessary already throughout $p\ge C_\eps\log n$, rather than only for
$p>(\log n)^2$.

\begin{corollary}[Uniform threshold improvement]\label{cor:uniform}
There are absolute constants $A,C>0$ and $n_0$ such that, for
$n\ge n_0$, $k\ge2$, $0<\eps<1/3$, and
\[
      p\ge A\frac{(\log n)^2}{\log\log n},
\]
the probability hypothesis of Corollary \ref{cor:eps} again implies
\eqref{eq:sharp}.
\end{corollary}

\begin{proof}
Write $L=\log n$ and $s=L/p$.  Then
\begin{equation}\label{eq:sbound}
 s\le \frac{\log L}{A L}.
\end{equation}
If $\eps\ge s$, then $\delta=(1+\eps)e^s-1<1/2$ for $n_0$ large enough,
and $\delta\le C_1\eps$.  The same monotonicity comparison used in
Corollary \ref{cor:eps} gives \eqref{eq:sharp}.

Suppose instead that $\eps<s$.  For $n_0$ large, \eqref{eq:sbound} gives
$s<1/16$, and
\[
 \delta=(1+\eps)e^s-1\le(1+s)e^s-1\le4s<1/2.
\]
Theorem \ref{thm:transfer} and monotonicity of
$t\mapsto t/\log(1/t)$ give
\[
 k\le C_2\frac{sL}{\log(1/s)}
 \le \frac{C_2}{A}
       \frac{\log L}{\log(AL/\log L)}
 \le \frac{2C_2}{A}.
\]
Choose the absolute constant $A$ so large that $2C_2/A<2$.  This contradicts
$k\ge2$.  Hence the probability hypothesis cannot hold in the case
$\eps<s$, and the remaining case already gives \eqref{eq:sharp}.
\end{proof}

\section{Verification and limitations}

The proof has one external dependency: the sharp randomized cube upper bound
\eqref{eq:cube}.  The transfer inequality \eqref{eq:distcompare} is exact and
pointwise in the subspace, so it loses neither probability nor dimension.
The estimates in the two corollaries only compare the increasing function
$t/\log(1/t)$ at arguments differing by absolute factors.

The result does \emph{not} settle the source question when
\[
 p\asymp\log n,
 \qquad
 \eps\ll\frac{\log n}{p}.
\]
In this wedge the global comparison error $n^{1/p}-1$ is larger than the
desired distortion.  A full solution appears to require a genuinely
subspace-sensitive analogue of the cube anti-concentration argument, rather
than a global norm comparison.

\section{Novelty check and human review}

A bounded search on 9 August 2026 covered the run's lightweight indexes, the
local arXiv source corpus, exact web/arXiv queries for the source title and the
logarithmic-$p$ question, and the OpenAlex citation graph of the source DOI
(24 citing records).  Lytova and Tikhomirov explicitly observed in 2017 that
the randomized Dvoretzky problem was not sharp for logarithmic $p$
\cite{LT18}; their later result resolves the Gaussian variance transition,
not this section-dimension question.  No later full resolution or this exact
two-parameter transfer statement was found.  This bounded search does not
prove novelty.

\medskip
\noindent\textbf{Recommendation: high-priority specialist verification.}
The proof is short.  A reviewer should confirm the exact normalization of
Tikhomirov's cube theorem and the use of its ``majority'' hypothesis, then
check the integer $k\ge2$ exclusion in Corollary \ref{cor:uniform}.

\begin{thebibliography}{9}

\bibitem{PVZ17}
G. Paouris, P. Valettas and J. Zinn,
\emph{Random version of Dvoretzky's theorem in $\ell_p^n$},
Stochastic Process. Appl. \textbf{127} (2017), 3187--3227;
arXiv:1510.07284.

\bibitem{Tik14}
K. E. Tikhomirov,
\emph{The randomized Dvoretzky's theorem in $\ell_\infty^n$ and the
$\chi$-distribution},
Geometric Aspects of Functional Analysis, Lecture Notes in Math. 2116,
Springer (2014), 455--463.

\bibitem{LT18}
A. Lytova and K. Tikhomirov,
\emph{The variance of the $\ell_p^n$-norm of the Gaussian vector, and
Dvoretzky's theorem},
Algebra i Analiz \textbf{30} (2018), 107--139;
arXiv:1705.05052.

\end{thebibliography}

\end{document}
